\SourceLeaf{015}\Article{10}
Although the relation between \(\lambda\) and \(\mu\) depends on the nature of the function \(\varphi x\), some general results can nevertheless be established. Namely, whatever this function may be, provided it is such that its value always decreases, or at least does not increase, as the absolute value of \(x\) increases, we shall certainly have
\begin{quote}
\(\lambda\) less than, or at least not greater than, \(\mu\BareRoot{3}\), whenever \(\mu\) is less than \(\tfrac23\);

\(\lambda\) not greater than \(\frac{2}{3\BareRoot{(1-\mu)}}\), whenever \(\mu\) is greater than \(\tfrac23\).
\end{quote}
For \(\mu=\tfrac23\) the two limits coincide: namely, \(\lambda\) cannot be greater than \(\BareRoot{\tfrac43}\).

To prove this remarkable theorem, let us denote by \(y\) the value of the integral \(\int\varphi z.\mathrm{d}z\) extended from \(z=-x\) to \(z=+x\), so that \(y\) will be the probability that an error is contained within the limits \(-x\) and \(+x\). Further, let us put
\[
x=\psi y,\quad \mathrm{d}\psi y=\psi'y.\mathrm{d}y,\quad \mathrm{d}\psi'y=\psi''y.\mathrm{d}y
\]
We shall therefore have \(\psi0=0\), and also
\[
\psi'y=\frac{1}{\varphi x+\varphi(-x)}
\]
so that, by hypothesis, \(\psi'y\) will always increase from \(y=0\) to \(y=1\), or at least will nowhere decrease; or, what is the same, the value of \(\psi''y\) will always be positive, or at least not negative. Further, we have \(\mathrm{d}.y\psi'y=\psi'y.\mathrm{d}y+y\psi''y.\mathrm{d}y\), and consequently
\[
y\psi'y-\psi y=\int y\psi''y.\mathrm{d}y
\]
the integration beginning at \(y=0\). The value of the expression \(y\psi'y-\psi y\) will therefore always be a positive quantity, or at least not a negative one, and consequently
\[
1-\frac{\psi y}{y\psi'y}
\]
will be a positive quantity less than unity. Let \(f\) be its value for \(y=\mu\), i.e. since we have \(\psi\mu=\lambda m\), let
\[
f=1-\frac{\lambda m}{\mu\psi'\mu}\quad\text{or}\quad\psi'\mu=\frac{\lambda m}{(1-f)\mu}
\]
With these preparations, let us consider the following function of \(y\)
\[
\frac{\lambda m}{(1-f)\mu}(y-\mu f)
\]
which we shall put \(=Fy\), and also \(\mathrm{d}Fy=F'y.\mathrm{d}y\). It is evident that
\SourceLeaf{016}
\[
\begin{aligned}
F\mu&=\lambda m=\psi\mu\\[3pt]
F'\mu&=\frac{\lambda m}{(1-f)\mu}=\psi'\mu
\end{aligned}
\]
Since, therefore, \(\psi'y\) increases continuously as \(y\) increases (or at least does not decrease, which is always to be understood), while \(F'y\) is constant, the difference \(\psi'y-F'y=\frac{\mathrm{d}(\psi y-Fy)}{\mathrm{d}y}\) will be positive for values of \(y\) greater than \(\mu\), and negative for smaller ones. From this it is easily concluded that \(\psi y-Fy\) is always a positive quantity, and consequently \(\psi y\) will always be absolutely greater, or at least not smaller, than \(Fy\), certainly as long as the value of \(Fy\) is positive, i.e. from \(y=\mu f\) to \(y=1\). Hence the value of the integral \(\int(Fy)^2\mathrm{d}y\) from \(y=\mu f\) to \(y=1\) will be smaller than the value of the integral \(\int(\psi y)^2\mathrm{d}y\) between the same limits, and therefore still more certainly smaller than the value of this integral from \(y=0\) to \(y=1\), which is \(=mm\). But the value of the former integral is found to be
\[
=\frac{\lambda\lambda mm(1-\mu f)^3}{3\mu\mu(1-f)^2}
\]
whence we conclude that \(\lambda\lambda\) is less than \(\frac{3\mu\mu(1-f)^2}{(1-\mu f)^3}\), where \(f\) is a quantity lying between \(0\) and \(1\). Now the value of the fraction \(\frac{3\mu\mu(1-f)^2}{(1-\mu f)^3}\), whose differential, if \(f\) is regarded as a variable quantity, is \(=\)
\[
-\frac{3\mu\mu(1-f)}{(1-\mu f)^4}.(2-3\mu+\mu f)\mathrm{d}f
\]
decreases continuously as \(f\) passes from the value \(0\) to the value \(1\), whenever \(\mu\) is less than \(\tfrac23\); consequently the greatest possible value will be that corresponding to \(f=0\), namely \(=3\mu\mu\), so that in this case \(\lambda\) is certainly less than or not greater than \(\mu\BareRoot{3}\). Q.~E.~P. On the other hand, whenever \(\mu\) is greater than \(\tfrac23\), the value of that fraction will be greatest for \(2-3\mu+\mu f=0\), i.e. for \(f=3-\tfrac2\mu\), whence it becomes \(=\frac4{9(1-\mu)}\), and therefore in this case \(\lambda\) is not greater than \(\frac2{3\BareRoot{(1-\mu)}}\). Q~E.~S.

Thus, for example, for \(\mu=\tfrac12\), \(\lambda\) certainly cannot be greater than \(\BareRoot{\tfrac34}\), i.e. the probable error cannot exceed the limit \(0{,}8660254m\), to which it was found equal in the first example of art. 9. Further, it is easily concluded from our theorem that \(\mu\) is not less than \(\lambda\BareRoot{\tfrac13}\) as long as \(\lambda\) is less than \(\BareRoot{\tfrac43}\); on the other hand, \(\mu\) is not less than \(1-\frac4{9\lambda\lambda}\) for a value of \(\lambda\) greater than \(\BareRoot{\tfrac43}\).

\Article{11}
Since several problems to be treated below are also connected with the value of the integral \(\int x^4\varphi x.\mathrm{d}x\), it will be worthwhile to develop it for certain spe\SourceLeaf{017}cial cases. We shall denote the value of this integral extended from \(x=-\infty\) to \(x=+\infty\) by \(n^4\).

I. For \(\varphi x=\tfrac1{2a}\), insofar as \(x\) is contained between \(-a\) and \(+a\), we have \(n^4=\tfrac15a^4=\tfrac95m^4\).

II. In the second case of art. 6, where \(\varphi x=\frac{a\mp x}{aa}\), for values of \(x\) between \(0\) and \(\pm a\), we obtain \(n^4=\tfrac1{15}a^4=\tfrac{12}{5}m^4\).

III. In the third case, where
\[
\varphi x=\frac{e^{-\frac{xx}{hh}}}{a\BareRoot{\pi}}
\]
we find, by what is set out in the memoir cited above, \(n^4=\tfrac34h^4=3m^4\).

Moreover, it can be proved that the value of \(\frac{n^4}{m^4}\) is certainly not less than \(\tfrac95\), provided that the hypothesis of the preceding article holds.

\clearpage
